\documentclass[../main]{subfiles}
\begin{document}
\chapter{Resistencia de un Conductor}

\img{images/01/resistencia}{Resistencia de un conductor}
\info{
Contenido}{
Resistencia de un conductor, resistividad.
}

\actividad{ Resolver ejercicios de resistividad}
\begin{enumerate}
    \item Un conductor de cobre tiene una resistencia de 1$\Omega$. Si se triplica su longitud.  ¿Cuál será su resistencia en Ohms?
    \item Si el conductor del problema anterior tiene sección cuadrada y  se duplica su sección.  ¿Cuál será su resistencia?
    \item Si la longitud inicial del conductor del problema 1 es de $20m$ de largo.  ¿Qué longitud habrá que acortar del conductor para que disminuya la resistencia a $0.4\Omega$?
    \item Un conductor de sección circular y 40 metros de longitud, tiene un diámetro de 2mm. Otro conductor mide 30 metros de largo y tiene un diámetro de 1mm. Si en ambos se mide el mismo valor de resistencia $4\Omega$,  ¿Están fabricados del mismo material?   
    \item Calcular la resistividad para un conductor que posee $2\Omega$, tiene un diámetro de 3mm y una longitud de 10 metros.
    \item Calcular la resistividad para un conductor que posee $10\Omega$, tiene un radio de 3mm y una longitud de 40 metros.
    \item Un cable metálico parece ser buen conductor, sobre una longitud 5m y 1mm de diámetro se midió $0.2\Omega$,  ¿Qué resistencia tendrá un cable fabricado con el mismo material de 40m largo y 2mm radio?
    \item Se conecta un conductor a una batería de $9V$ y se mide con un amperímetro que el cortocircuito marca $4A$. Si el conductor se corta a la mitad, y se conecta una de las partes.  ¿Cuánto debería marcar el amperímetro?
    \item  ¿Cuál es la resistividad del conductor del problema anterior si tiene una longitud de 2m y una sección de $2mm^2$?
    \item Un conductor de cobre tiene una resistividad $0.0171 \Omega mm^2/m$. Si tiene una longitud de 100m y una sección de $4mm^2$. Calcular su resistencia.
    \item Calcular el valor de la resistencia en Ohms($\Omega$) para un conductor de cobre con resistividad igual a la del problema anterior, que tiene una longitud de 1000m y un diámetro de $5mm$.
    \item La resistividad para el cobre tiene dos representaciones $0.0171 \Omega \frac{mm^2}{m}$ y $1,71 x 10^{-8}\Omega m$ Si el hierro tiene una resistividad de $8,90 x 10^{-8}\Omega m$  ¿Cual será su valor en $\Omega \frac{mm^2}{m}$ ?
    \item Elaborar una tabla con resistividades en $\Omega \frac{mm^2}{m}$ y $\Omega m$ para los materiales cobre,hierro,plata, oro,estaño, platino, aluminio y grafito.
    \end{enumerate}

\centering
\fbox{

$R=\rho (\frac{l}{s}) \Rightarrow \rho = R (\frac{s}{l}) $
}
\vspace{10mm}

Siendo:

$\rho$:Resistividad.

$s:$Sección del conductor.

$l:$longitud del conductor.

\end{document}